\documentclass{article}

\usepackage[utf8]{inputenc}
\usepackage[T1]{fontenc}
\usepackage{helvet}
\usepackage{geometry}
\usepackage{xcolor}
\usepackage{eso-pic}
\usepackage{fancyhdr}
\usepackage{parskip}
\usepackage{microtype}
\usepackage[english, ukrainian]{babel}
\usepackage{url}
\usepackage{amsmath}
\usepackage{amssymb}
\usepackage{amsthm}
\usepackage{longtable}
\usepackage{booktabs}
\usepackage{hyperref}
\usepackage{titlesec}

\definecolor{nato}{HTML}{293757}
\definecolor{footergray}{gray}{0.65}

\geometry{a4paper,left=3cm,right=3cm,top=3cm,bottom=3cm}
\renewcommand{\familydefault}{\sfdefault}
\setcounter{secnumdepth}{3}

\titleformat{\section}
  {\color{nato}\Huge\bfseries}{\thesection.}{1em}{}
\titlespacing*{\section}{0pt}{30pt}{12pt}

\titleformat{\subsection}
  {\color{nato}\LARGE\bfseries}{\thesubsection.}{1em}{}
\titlespacing*{\subsection}{0pt}{20pt}{8pt}

\titleformat{\subsubsection}
  {\color{nato}\Large\bfseries}{\thesubsubsection.}{1em}{}
\titlespacing*{\subsubsection}{0pt}{10pt}{6pt}

\fancyhf{}
\fancyhead[R]{\small\color{footergray} 31 березня 2026}
\fancyfoot[L]{\small\color{footergray} Технічні вимоги ЄСІКС. I. Інфраструктурні сервіси}
\fancyfoot[R]{\small\color{footergray}\thepage}

\newcommand\HUGE[1]{\fontsize{40}{40}\selectfont #1}

\renewcommand{\headrulewidth}{0pt}
\renewcommand{\footrulewidth}{0.4pt}
\renewcommand{\footrule}{\color{footergray}\hrule width \textwidth height 0.4pt}

\begin{document}

\pagestyle{fancy}

\begin{titlepage}
  \AddToShipoutPictureBG*{\AtPageLowerLeft{\color{nato}\rule{\paperwidth}{\paperheight}}}
  \color{white}\thispagestyle{empty}\vspace*{4cm}\centering
  {\Huge  \textbf{Технічні вимоги ЄСІКС} \par} \vspace{2cm}
  {\HUGE  \textbf{I. Інфраструктурні сервіси} \par} \vspace{2.5cm}
  {\LARGE \textbf{I-1. Інфраструктура безпеки (X.509, EUDI)} \par}
  {\LARGE \textbf{I-2. Рольова модель і система доступів (ABAC, IAM) } \par}
  {\LARGE \textbf{I-3. Оркестрація процесів (BPMN)} \par}
  {\LARGE \textbf{I-4. Словникова підсистема (HL7)} \par}
  {\LARGE \textbf{I-5. Мета сервіси схеми даних (SCHEMA)} \par}
  {\LARGE \textbf{I-6. Конструктор форм і процесів (X-FORMS)} \par} \vspace{2.5cm}
  \vfill
  {\large 2026 \copyright\ Інформаційні Судові Системи \par}
\end{titlepage}

\newpage

\begin{center}
  {\Huge\color{nato}\bfseries Зміст\par}
\end{center}
\vspace{40pt}
\tableofcontents

\begin{abstract}
ТЕХНІЧНІ ВИМОГИ
НА СТВОРЕННЯ ЗАСОБУ ІНФОРМАТИЗАЦІЇ - КОМПОНЕНТІВ “КЕРУВАННЯ ДОСТУПОМ”, БАЗОВИХ, БІЗНЕС СЕРВІСІВ ТА ІНФРАСТРУКТУРИ
ЄДИНОЇ СУДОВОЇ ІНФОРМАЦІЙНО-КОМУНІКАЦІЙНОЇ СИСТЕМИ.
\end{abstract}

\newpage
\section*{Вступна частина}

\subsection*{Умовні скорочення та визначення}

Терміни та скорочення, що використовуються в документі:

\begin{longtable}{p{4cm}p{10cm}}
\toprule
\textbf{Термін} & \textbf{Значення} \\
\midrule
Автентифікація & Електронна процедура, яка дає змогу підтвердити електронну ідентифікацію фізичної, юридичної особи, інформаційної або інформаційно-комунікаційної системи та/або походження та цілісність електронних даних \\
Авторизація & Частина процедури визначення та надання прав доступу до функцій та/або інформаційних ресурсів для роботи в інформаційній системі, після автентифікації \\
Валідація & Встановлення об’єктивних і задокументованих доказів того, що визначені вимоги до спеціально призначеного застосування можна послідовно виконати \\
API & Автоматизовані програмні інтеграції, відомі також як “прикладний інтерфейс програмування – application programming interface” \\
БД & База Даних — Організована колекція даних, яка зберігається та управляється з використанням комп'ютерної системи; забезпечує зберігання, організацію та доступ до даних з метою ефективного використання та обробки \\
БП & Бізнес Процес — Сукупність дій необхідних для здійснення людиною або системою, в результаті яких відбувається досягнення бажаного результату \\
КЕП & Кваліфікований електронний підпис (X.509) \\
РМК & Реєстраційно-моніторингова картка \\
АРМ & Автоматизоване робоче місце користувача \\
HSM & Hardware Secured Module \\
ЄСІКС & Єдина судова інформаційно-комунікаційна система \\
IAM & Identity and Access Management \\
SSO & Single Sign-On \\
BPMN & Business Process Model and Notation (ISO 19510) \\
ABAC & Attribute-Based Access Control (ISO 29146) \\
PEP & Policy Enforcement Point (ISO 29146) \\
PIP & Policy Information Point (ISO 29146) \\
PAP & Policy Administration Point (ISO 29146) \\
\bottomrule
\end{longtable}

\newpage
\subsection*{Загальні положення}

Технічні вимоги на основі «Концепції ЄСІКС» від 2024 року \cite{esiks}.

\subsection*{Призначення та цілі впровадження}

З метою забезпечення узгодженої, безпечної та ефективної роботи ЄСІКС,
а також унеможливлення фрагментації підходів до управління доступом,
даними та інфраструктурою, в рамках ЄСІКС є необхідним створення
та впровадження централізованих спільних компонентів, зокрема:

\begin{enumerate}
\item \textbf{Інфраструктура безпеки} IAM: управління доступами, користувачами, організаціями, ролями, ідентифікацією, автентифікацією, авторизацією, фіксація подій доступу.
\item \textbf{Рольова модель доступів} ABAC: централізоване керування правами на основі атрибутів, точок доступу, правил доступу.
\item \textbf{Оркестрація процесів} BPMN: управління документообігом, задачами, переходами, таймаутами, умовами, логікою.
\item \textbf{Словникова підсистема} HL7: системоутворюючі довідники, похідні піддовідники, версіонування, ієрархії, імутабельність.
\item \textbf{Мета сервіси схеми даних та конструктор форм}: конструктор сутностей, форм, бізнес-процесів, словників, правил доступу, шаблонів профілів користувачів.
\end{enumerate}

Реалізація зазначених складових є передумовою побудови цілісної, масштабованої та керованої архітектури ЄСІКС, зниження витрат на розробку та супровід, забезпечення єдиних стандартів безпеки та користувацького досвіду.

\section{Інфраструктура безпеки}
Інфраструктура відкритого ключа PKI X.509 ДСТУ 4145:2002 передбачає наступний перелік сервісів: CA, CMP, LDAP, CMS, CRL, OCSP, TSP.

\subsection{Провайдери української криптографії}
В Україні існують наступні провайдери криптографії ДСТУ:
\begin{itemize}
\item ТОВ ``Інститут Інформаційних Технологій''
\item АТ КБ ``ПриватБанк'' (Open Source)
\item ТОВ ``Сайфер''
\item ТОВ ``Автор'' (BouncyCastle)
\end{itemize}

\newpage
\subsection{Вимоги до сутностей}

\subsubsection{Директорія (LDAP)}
Директорія підприємства повинна опиратися на ITU/IETF: 2849, 3296, 3671-3673, 3866, 4510-4518, 4522, 4525, 4526, 4929, 5480, X.519.

\subsubsection{Сертифікати (X.509)}
Сервер сертифікатів повинен опиратися на ДСТУ 4145:2002 криптографію в стандартних X.509 конвертах і опиратися на стандарти ITU/IETF: 3279, 5755, 7030, X.509.

\subsubsection{Користувачі (inetOrgPerson)}
Реквізитна частина Користувача повинна бути зафіксована технічним завданням.

\subsubsection{Організації (organization)}
Реквізитна частина Організації повинна бути зафіксована технічним завданням.

\subsubsection{Криптографічне повідомлення (CMS X.894)}
Криптографічне повідомлення CMS повинно відповідати стандарту CMS X.894 і підтримувати криптографію ДСТУ 4145:2002.

\subsubsection{Кваліфікований електронний підпис (ДСТУ 4145)}
КЕП повинен відповідати стандарту CMS X.894 і підтримувати криптографію ДСТУ 4145:2002.

\subsubsection{Відкликання і перевірка (OCSP)}
У складі інфраструктури PKI X.509 повинен бути сервер перевірки валідності сертифікатів або управління CRL списками.

\subsubsection{Мітка часу (TSP)}
У складі інфраструктури PKI X.509 повинен бути сервер позначки часу для часової аутентифікації.

\subsubsection{Ключі шифрування та сервісні креденціали}
Для використання у складі апартного HSM або програмного комплекту з захищеним сховищем.

\newpage
\subsection{Вимоги до процесів}
Над сутностями системи інфраструктура безпеки повинна експонувати наступні процеси.

\subsubsection{Видача сертифікату (enroll)}
Сервер сертифікатів повинен генерувати сертифікати по PKCS-10 CSR протоколу через TCP/DER або HTTP/DER версії CMP протоколів.

\subsubsection{Автентифікація сертифікатом (check)}
За допомогою клієнтського сертифікату і сервера часу здійснюється підпис автентифікаційного повідомлення і реєстрація
в системі сесій і пристроїв (SSO). Кожен користувач повинен мати змогу побачити свої сесії через систему IAM.

\subsubsection{Відкликання сертифікату (revoke)}
Система повинна містити процедуру відкликання сертифікату.

\subsubsection{Оновлення сертифікату (renew)}
Продовження терміну дії сертифікату.

\subsubsection{Підпис і верифікація (sign/verify)}
На кожному пристрої, який може здійснювати SSO логін повинна бути імплементова бібліотека криптографічних операцій ДСТУ 4145:2002 або ECDSA.

\subsubsection{Шифрування і розшифрування (encrypt/decrypt)}
На кожному пристрої, який може здійснювати шифрування і розшифрування повинна бути імплементова бібліотека
криптографічних операцій ДСТУ 7624:2014, 7564:2014 або AES-GCM-256, AES-CBC-256, XSalsa20-Poly1305, XChaCha20-Poly1305.

\subsubsection{Управління секретами (store/retrieve)}
Управління ключами шифрування і сервісними креденціалами.

\subsection{Вимоги до серіалізації і сховища}
Сертифікати X.509 описуються ASN.1 X.680 мовою і кодуються X.690 DER серіалізатором. Для токенів Європейської Системи EUDI і
їх нод eIDAS використовується CBOR (RFC 8742, 8949, 7049) серіалізація \cite{cbor}.
Сховище повинно бути побудоване на низьколатентних бібліотеках з єдиних простором ключів (LSM-дерева) і кешуванням в пам'яті
оскільки сервери сертифікатів є вузьким місцем для всіх залежних сервісів.

\addtocontents{toc}{\protect\newpage}
\newpage
\section{Рольова модель і система доступів}

\subsection{Вимоги до рольової моделі ABAC}

ЄСІКС реалізує систему керування доступом на основі атрибутів (ABAC) згідно з міжнародним
стандартом ISO/IEC 29146:2016 (Information technology — Security techniques — A framework for
access management) та рекомендаціями NIST SP 800-162.

\subsubsection{Об'єкт доступу}

Серед об'єктів системи ABAC повинен бути наступний мінімум:

\begin{itemize}
\item Співробітник
\item Процес
\item Документ
\end{itemize}

\subsubsection{Суб'єкт доступу}

Суб'єктами системи можуть виступати будь-які автентифіковані PKI X.509 користувачі.

\subsubsection{Правило доступу}

Згідно ISO 29146 правила можуть мати типи, групуватися в політики операціями Булевої алгебри,
поєднують суб'єкт і об'єкт авторизації. Правило доступу визначає умови позитивної функції
на реквізитній частині об'єкту верифікації і суб'єкту верифікації в точці підключення.

\subsection{Вимоги до сутностей ABAC}

\subsubsection{Політики доступу}
Політики доступу представляють набір правил доступу згруповані за певною ознакою.

\subsection{Вимоги до процесів ABAC}

\subsubsection{Життєвий цикл Політики доступу}

Створення, редагування, реєстрація і видалення політики доступу.

\newpage
\subsubsection{Точки доступу}

Перелік обов'язкових точок доступу в ЄСІКС згідно ISO/IEC 29146 від 2016-06-01 визначених як PEP для BPMN.

\begin{itemize}
\item BLOCK Блокування користувача
\item SEARCH Пошук документів
\item VIEW Перегляд папки
\item CREATE Створення документу
\item EDIT Редагування форми
\item DELETE Видалення процесу
\item NEXT Наступний крок процесу
\item REJECT Відхилення
\item SIGN Підпис
\item REGISTER Реєстрація документу
\item CANCEL Відхилення
\item SEND Відправка
\item RECEIVE Отримання
\item STATUS Модифікація статусу документу
\end{itemize}

\subsection{Вимоги до сутностей IAM}

ЄСІКС повинна мати систему керування доступом згідно з міжнародним
стандартом ISO/IEC 27001 Information security management systems,
ISO/IEC 24760-1:2019 IT security and privacy — A framework for identity management,
ISO/IEC 27018:2019.

\subsubsection{Адміністратор організацій}
Адміністратор організацій здійснює облік підприємств судової системи.

\subsubsection{Адміністратор структурних підрозділів}
Адміністратор користувачів здійснює управління структурними підрозділами певного підприємства судової системи.

\subsubsection{Адміністратор співробітників}
Адміністратор користувачів здійснює управління співробітниками конкретного підприємства судової системи.

\subsubsection{Адміністратор користувачів}
Адміністратор користувачів здійснює управління користувачів будь-якого підприємства судової системи.

\newpage
\subsection{Вимоги до процесів IAM}

\subsubsection{Життєвий цикл Адміністратора}
Створення, редагування, деактивація адміністраторів системи.

\subsubsection{Життєвий цикл Організації}
Створення, редагування, деактивація організацій системи.

\subsubsection{Життєвий цикл Користувача}
Створення, редагування, деактивація користувачів системи.

\subsubsection{Життєвий цикл Сесії}
Кожен користувач системи повинен мати змогу управляти своїми (SSO X.509) сесіями.

\subsubsection{Життєвий цикл Ролі}
Ролі системи IAM пов'язані з точками доступу ABAC PEP.

\subsection{Вимоги до рольової моделі IAM}

\subsubsection{Точки доступу}
Перелік обов'язкових точок доступу в ЄСІКС згідно ISO/IEC 29146 від 2016-06-01 визначених як PEP для IAM.

\begin{itemize}
\item VIEW Переглянути реєстр користувачів/організацій
\item CARD Переглянути картку користувача/організації
\item CREATE Створити користувача/організацію
\item EDIT Редагувати картку користувача/організацію
\item DELETE Деактивувати користувача/організацію
\end{itemize}

\subsection{Вимоги до сховища}
Сховище повинно бути побудоване на низьколатентних бібліотеках з єдиних простором ключів (LSM-дерева) і кешуванням в пам'яті
оскільки це міце є вузьким місцем для всіх залежних сервісів.

\subsection{Вимоги до серіалізації і валідації}

Серіалізація повинна відбуватися в JSON, а валідація об'єктів здійснюватися черех JSON Schema (draft-07).

\newpage
\section{Оркестрація процесів}

\subsection{Вимоги до сутностей BPMN}
ЄСІКС реалізує систему управління діло-процесами згідно з міжнародним
стандартом ISO/IEC 19510 Information technology — Object Management Group Business Process Model and Notation.

\subsubsection{Документ}
Об'єк похідний від базового документу. Повна таксономія документів і їх реквізитні частини визначаються технічним завданням.

\subsubsection{Процес}
Процес визначає абстракцію над BPMN XML визначенням процесу, яка має здатність інстанціюватися.

\subsubsection{Монітор}
Абстракція, що представляє ієрархічну групу процесів для утворення реєстраційно-моніторингової картки документообігу.

\subsubsection{Задача користувача (АРМ)}
Вузол графа процесу, що представляє собою юридичне підтвердження людського втручання.

\subsubsection{Системна задача}
Вузол графа процесу, що представляє собою автоматичну задачу.

\subsubsection{Перехід}
Ребро графа процесу, яке містить функцію на основній мові програмування імплементації або скриптовій мові,
що має доступ до контекста процесу (всіх його документів на поточній стадії, і повну історію з документами на кожній стадії).

\newpage

\subsubsection{Подія}
Асинхронна зовнішня подія, яка залишає слід в історії.

\subsubsection{Таймаут}
Подія, яка автоматично виконує певну функцію процесу.

\subsubsection{Асинхронне повідомлення}
Асинхронна відправка довільного повідомлення в інстанцію процесу.

\subsection{Вимоги до процесів}

\subsubsection{Створення}
Реєстрація нової інстанції бізнес-процесу на основі завантаженого XML-визначення за стандартом ISO 19510.

\subsubsection{Зупинення}
Механізм призупинення виконання поточних задач інстанції процесу зі збереженням її повного стану.

\subsubsection{Відновлення}
Процедура активації раніше зупиненої інстанції процесу з відновленням виконання з останньої точки збереження.

\subsubsection{Архівація}
Перенесення даних завершених або неактивних інстанцій процесів до довгострокового сховища для звільнення оперативних ресурсів.

\subsubsection{Зміна контексту процесу}
Динамічна модифікація внутрішніх даних та метаданих інстанції процесу під час її виконання.

\subsubsection{Перехід на стадію}
Управління переміщенням токена виконання між вузлами графа процесу відповідно до логіки переходів.

\subsection{Вимоги до сховища}
Сховище повинно забезпечувати атомарну персистентність стану процесів, низьку затримку доступу та підтримку транзакцій.

\addtocontents{toc}{\protect\newpage}
\newpage
\section{Словникова підсистема}
ЄСІКС реалізує систему управління похідними версіоновами ієрархічними словниками згідно з міжнародним
стандартом ISO/IEC HL7.

\subsection{Загальні відомості}

\begin{enumerate}
\item Можливість проведення модернізації довідників без контрактів розробки
\item Підтримка версіонування словників
\item Підтримка ієрархічних словників
\item Розробка ергономічної адміністративної панелі оператора
\item Сумісність з існуючими записами ІСС і збереження історичних даних
\item Відповідність міжнародним стандартам ISO
\item Імутабельність даних (ніщо не видаляється з системи)
\item Максимальна повнота (ієрархії, трансляція кодів, валідація)
\item Публічний сервіс для пошуку, перегляду, порівняння, експорту
\item Приватний сервіс адміністраторів і операторів ДСА для редагування і публікації, імпорту
\end{enumerate}

\subsection{Вимоги до сутностей HL7}

\subsubsection{Системноутворюючі довідники}
Сутність Code System використовується для системоутворюючих довідників,
оригінальних довідників які містять повну множину значень, юридичні підписані
довідники без права модифікації, імпортовані з інших систем, унормовані національні переклади, тощо.

\subsubsection{Похідні довідники}
Сутність Value Set використовується для операційних похідних довідників,
що відображаються в медичних інформаційних системах. Вони часто містять
підмножини значень існуючих довідників, можуть бути складені з існуючих довідників,
а також містити повні переліки значень CodeSystem, але мати іншу атрибутику і призначення.

\subsubsection{Система імен}
Ідентифікація словників повинна здійснюватися через унікальні OID або URI для забезпечення глобальної унікальності в ЄСІКС.

\subsubsection{Відображення довідників}
Використання сутності ConceptMap для встановлення відповідності між кодами різних довідникових систем.

\subsubsection{Налаштування довідникової системи}
Глобальна конфігурація параметрів валідації, кешування часто використовуваних значень та правил версіонування.

\newpage
\subsection{Вимоги до процесів}

\subsubsection{Створення попереднього словника}
Ініціалізація нової чернетки словника для наповнення концептами та атрибутами.

\subsubsection{Редагування попереднього словника}
Внесення змін до структури, назв або кодів у межах непублікованої версії словника.

\subsubsection{Публікація словника}
Надання словнику статусу "Active" та фіксація його імутабельної версії для використання в системі.

\subsubsection{Деактивація словника}
Припинення використання словника для нових записів зі збереженням доступу для перегляду історичних даних.

\subsubsection{Архівування словника}
Переміщення застарілих версій словників до архівного сховища згідно з регламентами зберігання даних.

\subsection{Вимоги до сховища}
Сховище словників повинно підтримувати ієрархічну структуру, імутабельність версій та швидкий пошук за кодами.

\newpage
\section{Мета сервіси та конструктор}
\subsection{Реєстр сутностей системи (API)}
Централізований облік усіх програмних інтерфейсів та сутностей системи з детальним описом їх реквізитної частини.

\subsection{Біблітека валідації (JSON Schema)}
Набір схем у форматі JSON Schema для автоматичної перевірки цілісності та коректності вхідних даних.

\subsection{Конструктор сутностей та форм}
Інструментарій для візуального створення структур даних та відповідних користувацьких інтерфейсів (X-Forms).

\subsection{Конструктор бізнес процесів}
Графічний редактор для моделювання логіки документообігу та автоматизованих задач у форматі BPMN.

\subsection{Конструктор словників}
Інтерфейс для керування таксономіями, ієрархіями та наповненням довідникової підсистеми ЄСІКС.

\subsection{Конструктор правил доступу}
Засіб конфігурації політик ABAC на основі атрибутів суб'єктів та об'єктів доступу.

\subsection{Конструктор шаблонів профілів користувачів}
Інструмент для визначення набору атрибутів та ролей, що застосовуються до категорій користувачів системи.

\newpage
\section*{Висновок}

Перша частина сервісів ЄСІКС «Інфраструктурні сервіси» закладає фундамент єдиного інформаційного простору судової системи України.
Впровадження централізованої інфраструктури безпеки (X.509/EUDI), гнучкої рольової моделі на основі атрибутів (ABAC),
оркестрації бізнес-процесів за стандартом BPMN та унормованої словникової підсистеми (HL7) забезпечує цілісність,
масштабованість та високий рівень захисту даних.

\begin{thebibliography}{9}

\bibitem{ecoz:rehab} Міністерство охорони здоров’я України. \newblock Технічні вимоги Реабілітація 2.0. \newblock Версія 08.08.2024.
\bibitem{dstu:3973} ДСТУ 3973-2000. \newblock Система розроблення та постачання продукції на виробництво. Правила виконання науково-дослідних робіт. Загальні положення.
\bibitem{dstu:3974} ДСТУ 3974-2000. \newblock Система розроблення та постачання продукції на виробництво. Правила виконання науково-конструкторських робіт. Загальні положення.
\bibitem{dstu:42010} ДСТУ 42010:2018. \newblock Інженерія систем і програмних засобів. Опис архітектури.
\bibitem{law:info-protection} Закон України «Про захист інформації в інформаційно-телекомунікаційних системах».
\bibitem{law:personal-data} Закон України «Про захист персональних даних».
\bibitem{law:cybersecurity} Закон України «Про основні засади забезпечення кібербезпеки України».
\bibitem{owasp:top10} OWASP Foundation. \newblock OWASP Top 10: The Ten Most Critical Web Application Security Risks. \newblock \url{https://owasp.org/www-project-top-ten/}.
\bibitem{nist:sp800-57} NIST Special Publication 800-57 Part 1 Revision 5. \newblock Recommendation for Key Management.
\bibitem{iso:42010} ISO/IEC/IEEE 42010:2022. \newblock Systems and software engineering — Architecture description.
\bibitem{zachman} Zachman, J.A. \newblock A Framework for Information Systems Architecture. \newblock IBM Systems Journal, 1987.
\bibitem{esiks} Концепція ЄСІКС. \url{https://court.gov.ua/storage/portal/dsa/news/Kontseptsia%20ESIKS.pdf}. 2024
\bibitem{cbor} CBOR/COSE encoding/protocol. \url{https://tonpa.guru/stream/2024/2024-11-20%20CBOR%20COSE.htm}. 2025

\end{thebibliography}

\end{document}
